\documentclass[a4paper,12pt,oneside]{book} % nie: report!


% pakiety
\usepackage{polski} % lepiej to zamiast babel!
\usepackage[utf8]{inputenc} % w razie kłopotów spróbować: \usepackage[utf8x]{inputenc}
\usepackage{fancyhdr} % nagłówki i stopki
\usepackage{indentfirst} % WAŻNE, MA BYĆ!
\usepackage[pdftex]{graphicx} % to do wstawiania rysunków
\usepackage{amsmath} % to do dodatkowych symboli, przydatne
\usepackage[pdftex,
            left=1in,right=1in,
            top=1in,bottom=1in]{geometry} % marginsy
\usepackage{amssymb} % to też do dodatkowych symboli, też przydatne
\usepackage{pdfpages}
\usepackage{lipsum}
\usepackage{multirow}
\usepackage{listings}
\usepackage{caption}
\DeclareCaptionType{code}[Listing][Spis listingów] 

% definicje nagłówków i stopek
\pagestyle{fancy}
\renewcommand{\chaptermark}[1]{\markboth{#1}{}}
\renewcommand{\sectionmark}[1]{\markright{\thesection\ #1}}
\fancyhf{}
\fancyhead[LE,RO]{\footnotesize\bfseries\thepage}
\fancyhead[LO]{\footnotesize\rightmark}
\fancyhead[RE]{\footnotesize\leftmark}
\renewcommand{\headrulewidth}{0.5pt}
\renewcommand{\footrulewidth}{0pt}
\addtolength{\headheight}{1.5pt}
\fancypagestyle{plain}{\fancyhead{}\cfoot{\footnotesize\bfseries\thepage}\renewcommand{\headrulewidth}{0pt}}


% interlinia
\linespread{1.25}


% treść
\begin{document}
\sloppy
\thispagestyle{empty}
\includepdf{stronatytulowa}
\newpage{}

\thispagestyle{empty}
\newpage{}

\tableofcontents{}

\chapter*{Wstęp}
\addcontentsline{toc}{chapter}{Wstęp}
\label{Wstep}
%Mega ogólne sprawy dotyczące tego co chce zrobić
Sieci neuronowe mają za zadanie naśladowanie zachowań sieci neuronów znajdujących się w mózgu człowieka. Zostały stworzone do rozwiązywania zadań trudnych lub prawie niemożliwych do opisania za pomocą reguł, wyrażeń logicznych i innych narzędzi programistycznych. Wraz z rozwojem sieci neuronowych powstało wiele wariantów, które ze względu na swoją budowę lepiej lub gorzej sprawdzają się w różnych problemach. W przypadku rozpoznawania obrazów w postaci dwu-wymiarowej macierzy dla danych monochromatycznych lub trój-wymiarowej macierzy dla zdjęć kolorowych jednym z najlepszych wyborów będą sieci konwolucyjne. Sieci te rozpoznają wzorce, poczynając od linii horyzontalnych i wertykalnych, a w dalszych warstwach kończąc na skomplikowanych strukturach. Budowa i działanie sieci konwolucyjnych daje wielki potencjał do klasyfikacji obrazów. Ta praca przedstawi przykład takiej klasyfikacji wieloklasowej z użyciem sieci konwolucyjnych na przykładzie alfabetu Amerykańskiego Języka Migowego. 

\chapter*{Cel i zakres pracy}
\addcontentsline{toc}{chapter}{Cel i zakres pracy}
\label{Cel i zakres pracy}
%celem pracy jest napisanie programu bla bla
Celem pracy jest stworzenie programu wyposażonego w wytrenowany model do rozpoznawania obrazu, który w czasie rzeczywistym używając kamery internetowej będzie w stanie odczytać i wyświetlić na ekranie transkrypcję znaków języka migowego pokazywanych przez osobę znajdującą się w polu widzenia kamery. Dodatkowo do pracy będą składać się: utworzenie zbioru danych składającego się z około 50000 zdjęć zawierających wszystkie litery alfabetu ASL, utworzenie modelu z warstw konwolucyjnych i gęstych oraz wytrenowanie modelu i tuning parametrów. W części teoretycznej pracy zostanie przybliżony temat języka migowego American Sign Language, jak również temat konwolucyjnych sieci neuronowych i sposobu działania modeli.

\chapter{Język migowy American Sign Language}

\section{Historia}
\section{Populacja}
\section{Dystrybucja geograficzna}
\section{Składnia języka}


\chapter{Konwolucyjne sieci neuronowe i rozpoznawanie obrazów}

\section{Sekcja C}
\section{Sekcja D}

\chapter{Implementacja aplikacji do rozpoznawania języka migowego ASL}

\chapter{Podsumowanie}


\listoftables{} % jeśli są tabele
\addcontentsline{toc}{chapter}{Spis tabel}

\listoffigures{} % jeśli są tabele
\addcontentsline{toc}{chapter}{Spis rysunków}

\listofcodes
\addcontentsline{toc}{chapter}{Spis listingów}

\addcontentsline{toc}{chapter}{Bibliografia}
\bibliographystyle{ieeetr}
\bibliography{bibliography/cite}


\end{document}
